\subsection{Descent from completely split primes}
\label{subsec:split-prime-descent}

\begin{proof}[Proof of Theorem~\ref{thm:split-prime-descent}]
Take $x=(x_v)\in X(\A_\Q)_\bullet^{\mathrm{f\text{-}ab}}$.
Functoriality gives $f(x)\in A(\A_\Q)_\bullet^{\mathrm{f\text{-}ab}}$.
For every rational prime $\ell$, its componentwise $\ell$-adic
Kummer image comes from $\Sel(T_\ell A)$: pull back the
multiplication-by-$\ell^r$ torsors and take a compatible inverse
limit of the finite Selmer sets. The exact sequence is
\[
0\longrightarrow A(\Q)^{\wedge_\ell}
\longrightarrow\Sel(T_\ell A)
\longrightarrow T_\ell\Sha(\Q,A)\longrightarrow0.
\]
At each completely split prime satisfying the hypothesis, this
Selmer class is induced by some $a_\ell\in A(\Q)[\ell^\infty]$.
The difference $f(x_\ell)-a_\ell$ has zero $\ell$-adic Kummer
class in $A(\Q_\ell)$. A compact abelian $\ell$-adic Lie group
has a finite-index open subgroup isomorphic to $\Z_\ell^d$;
therefore the kernel of its pro-$\ell$ completion is precisely
its finite prime-to-$\ell$ torsion. Consequently $f(x_\ell)$
is torsion. This assertion does not require good reduction.

By Manin--Mumford, in the form of
\citet[Theorem~1.1$^*$, p.~151]{PilaZannier2008}, the hypothesis
on the image curve implies that
\[
f(X)(\bar\Q)\cap A(\bar\Q)_{\rm tors}
\]
is finite and Galois stable. Its inverse image is a finite reduced
closed subscheme $Z\subset X$ defined over $\Q$, since $f$ is
nonconstant. Thus $x_\ell\in Z(\Q_\ell)$ at all the indicated
split primes.

Finite abelian descent survives base change by
\citet[Proposition~5.16, p.~370]{Stoll2007}; the abelian assertion
is explicitly retained in \citet[\S2(4)]{StollErrata2017}. Hence
\[
x_L\in X(\A_L)_\bullet^{\mathrm{f\text{-}ab}}.
\]
At every place $v\mid\ell$ above one of the indicated completely
split primes, $(x_L)_v\in Z(L_v)$. These places have Dirichlet
density one among the finite places of $L$. Indeed, outside
ramification, a prime ideal has residue degree one exactly when
its rational prime splits completely, since $L/\Q$ is Galois.
The sum of $\operatorname{Norm}(v)^{-s}$ over residue degrees at
least two stays bounded as $s\to1^+$, whereas the sum over all
prime ideals diverges logarithmically. Finitely many exceptions
do not change the density.

A nonconstant map from $X$ to an abelian variety implies
$g(X)\ge1$. Therefore
\citet[Theorem~8.2, pp.~383--384]{Stoll2007}, in its
density-one finite-support form, gives $x_L=P$ for some
$P\in Z(L)$. The base-changed adelic point is Galois invariant.
The diagonal map $X(L)\to X(\A_L)_\bullet$ is injective, as
one sees in any finite completion, so $P\in Z(\Q)$.

For each finite rational prime $p$, choose $v\mid p$.
Equality over $L_v$ and the injectivity
$X(\Q_p)\to X(L_v)$ give $x_p=P$. Base change to a field
with complex places can discard the real component. Apply
Stoll's Theorem~8.2 a second time over $\Q$, to the singleton
$\{P\}$ and all finite places. This restores equality in the
entire modified adelic space. The validity of using finite
places is also explicit in \citet[\S3]{StollErrata2017}.
The converse follows because rational points survive descent.
\end{proof}

\begin{remark}[Geometric and arithmetic hypotheses]
\label{rem:split-prime-scope}
The geometric hypothesis holds if $\dim A\ge2$ and the
homomorphism $\Jac(X)\to A$ induced by $f$ is surjective:
differences of an elliptic translate generate only that
elliptic subvariety. Genus at least two of $X$ alone does not
exclude an elliptic image. Over any number field, the same
initial argument followed directly by Stoll's Theorem~8.2
works if the places whose residue characteristic $\ell$
satisfies $T_\ell\Sha(k,A)=0$ have density one. The completely
split criterion avoids imposing that density-one condition
over $\Q$ itself. Neither formulation proves the required
arithmetic vanishing for a new abelian variety.
\end{remark}
